%! TEX root = ./main.tex

\section{\texorpdfstring{Proof of the main result: case $\beta \in \R$}{}}



For $\beta \in (1,2) \backslash \Q$, one cannot cast (\ref{eq:greedy map chunks as dynamical system}) into an effective algorithm, because it uses comparisons with $\beta$ in its computations, which might require infinite precision. Here, we instead define a dynamical system that only manipulates finite-precision approximations of $\beta$. Inspired from the considerations which applied for the simpler case $\beta \in \Q \cap (1,2)$, we introduce a dynamical system which generates chunks of a $\beta$-expansion of a given $s \in [0,1]$ by reading increasingly long prefixes of both the greedy binary expansion of $s$ and $\beta - 1$. For computable $\beta \in (1,2)$, the generated $\beta$-expansion will then be proven to satisfy (\ref{eq:inequality to be proven}). Let $s \in [0,1]$ and $N\in \N$. When presented with increasingly long prefixes of $\gf_2(s)$ of respective lengths $\Sigma_i \in \N$, and increasingly long prefixes of $\gf_2(\beta)$ of respective lengths $\Lambda_i \in \N$, the dynamical system converts $\langle \gf_2(s)_{1: \Sigma_i}, \gf_2(\betaun)_{1: \Lambda_i}\rangle$ to $\xf_{1: Ni}$, for $i \in \N$, where $\xf \in \Sigma_\beta(s)$. We will choose the $\Sigma_i$ and $\Lambda_i$ to be such that $\limi{i} \frac{\Sigma_i}{Ni} = \log_2(\beta)$ and $\limi{i} \frac{\Lambda_i}{Ni} = 1$, which, in turn, leads to (\ref{eq:inequality to be proven}) for $\beta$ computable. Specifically, we shall see that the dynamical system is realized by an effective algorithm. Figure \ref{fig:principle conversion greedy binary to beta exp beta r} illustrates this process. 
\begin{figure}[H]
	\centering
	\tikzmath{\x=3;}
	\tikzmath{\y1=0.25;}
	\tikzmath{\y2=0.5;}
	\tikzmath{\S=0.93;}
	\begin{tikzpicture}[scale = 0.8]
		\node[anchor = west] at (9.35,0.1){{\small $\gf_2(\betaun) = 01001|0001|111|1001|000|1101|\ldots$}};
		\node[anchor = west] at (0,0.1){{\small $\gf_2(s) = 01001|0001|111|1001|000|1101|\ldots$}};
		\node[anchor = west] at (5.7,-\x){{\small$\xf = \underset{N}{000100}|\underset{N}{101000}|\underset{N}{001000}|101000|\ldots$}};
		\draw[thick, ->] (2.8,-1.7) -- (2.9*\S+4.1, -\x + \y2-0.1);
		\draw[thick, ->] (4*\S,-1) -- (4.6*\S+4.5, -\x + \y2);
		\draw[thick, ->] (5*\S,-0.4) -- (6.1*\S+4.5, -\x + \y2);
		\draw [decorate,decoration={brace,amplitude=5pt,mirror}]
  (1.9,-1.5) -- (3.1,-1.5) node[midway, yshift = -0.9em]{\scriptsize$\Sigma_1$};
  \draw [decorate,decoration={brace,amplitude=5pt,mirror}]
  (1.9,-0.8) -- (4.2,-0.8) node[midway, yshift = -0.9em]{\scriptsize$\Sigma_2$};
  \draw [decorate,decoration={brace,amplitude=5pt,mirror}]
  (1.9,-0.22) -- (5,-0.22) node[midway, yshift = -0.9em]{\scriptsize$\Sigma_3$};
  
  \draw[thick, ->] (12,-0.34) -- (2.9*\S+5, -\x + \y2);
		\draw[thick, ->] (12,-0.92) -- (4.6*\S+5, -\x + \y2);
		\draw[thick, ->] (5*\S+7.6,-1.64) -- (6.1*\S+5.5, -\x + \y2);
		\draw [decorate,decoration={brace,amplitude=5pt,mirror}]
  (11.9,-0.22) -- (13.1,-0.22) node[midway, yshift = -0.9em]{\scriptsize$\Lambda_1$};
  \draw [decorate,decoration={brace,amplitude=5pt,mirror}]
  (11.9,-0.8) -- (14.2,-0.8) node[midway, yshift = -0.9em]{\scriptsize$\Lambda_2$};
  \draw [decorate,decoration={brace,amplitude=5pt,mirror}]
  (11.9,-1.5) -- (15,-1.5) node[midway, yshift = -0.9em]{\scriptsize$\Lambda_3$};
	\end{tikzpicture}
	\caption{Conversion of the greedy binary expansions $\gf_2(s)$ and $\gf_2(\beta \hspace{-0.5mm} - \hspace{-0.5mm} 1)$ to a $\beta$-expansion $\xf$ of similar algorithmic complexity.}
	\label{fig:principle conversion greedy binary to beta exp beta r}
\end{figure}

The chunks of the $\beta$-expansion $\xf$ are generated by the following iterative process. Suppose that at step $i \in \N_0$, the first $i$ chunks of $\xf$ of size $N$ have been generated from has been generated from $\gf_2(s)_{1:\Sigma_i}$ and $\gf_2(\betaun)_{1:\Lambda_i}$ Now, we read the next chunk of $\gf_2(s)$, of size $\Sigma_{i+1}-\Sigma_i$, and the next chunk of $\gf_2(\betaun)$, of size $\Lambda_{i+1}-\Lambda_i$. We define 
\begin{equation}\label{eq:beta i to generate beta expansion beta r}
	\tag{\text{$\mathcal B$}}
	\beta_{i+1} := 1+\delta_2(\gf_2(\betaun)_{1:\Lambda_{i+1}}),
\end{equation}
and we generate a chunk $\xf$ of size $N$ by choosing the lexicographically maximal binary sequence $x \in \{0,1\}^N$ that satisfies 
\begin{equation}\label{eq:reducing the error between x beta and greedy binary beta r}
	\delta_2(\gf_2(s)_{1:\Sigma_{i+1}}) - \beta_{i+1}^{-N(i+1)} \leq \delta_{\beta_{i+1}}(\xf_{1:Ni}x) \leq \delta_2(\gf_2(s)_{1:\Sigma_{i+1}}),
\end{equation}
and take $\xf_{1:N(i+1)} := \xf_{1:Ni} x$. Using Lemma \ref{lem:factorisation of beta expansion through concatenation}, one can rewrite (\ref{eq:reducing the error between x beta and greedy binary beta r}) as
\begin{equation}\label{eq:first step towards x greedy exp of R i beta r}
	\delta_2(\gf_2(s)_{1:\Sigma_{i+1}}) - \frac{\beta_{i+1}^{-N(i+1)}}{\beta_{i+1} - 1} \leq \delta_{\beta_{i+1}}(\xf_{1:Ni}) + \beta^{-Ni} \delta_{\beta_{i+1}}(x) \leq \delta_2(\gf_2(s)_{1:\Sigma_{i+1}}).
\end{equation}
By defining $R_0 = 0$, $\Sigma_0 = 0$,
\begin{equation}\label{eq:R i definition beta r}
		R_i := \beta_{i+1}^{Ni}\left(\delta_2(\gf_2(s)_{1:\Sigma_{i}}) - \delta_{\beta_{i}}(\xf_{1:Ni})\right), \ \text{for} \ i \in \N,
	\tag{\text{$\mathcal{R'}$}}
\end{equation}
\begin{equation}\label{eq:s i to generate beta expansion beta r}
	s_i := \frac{\beta_{i+1}^{Ni}}{2^{\Sigma_i}}\delta_2(\gf_2(s)_{\Sigma_i+1:\Sigma_{i+1}}), \ \text{for} \ i \in \N_0,
	\tag{\text{$\mathcal{S'}$}}
\end{equation}
and
\begin{equation}\label{eq:epsilon i to generate beta expansion beta r}
	\varepsilon_i := \beta_{i+1}^{Ni} \left(\delta_{\beta_{i}}(\xf_{1:Ni}) - \delta_{\beta_{i+1}}(\xf_{1:Ni})\right), \ \text{for} \ i \in \N_0, \tag{\text{$\varepsilon$}}
\end{equation}
(\ref{eq:first step towards x greedy exp of R i}) is equivalent to 
\begin{equation}\label{eq:greedy exp for x beta r}
	R_i + s_i + \varepsilon_i - \frac{\beta_{i+1}^{-N}}{\beta_{i+1}-1} \leq \delta_{\beta_{i+1}}(x) \leq R_i + s_i + \varepsilon_i.
\end{equation}
$x \in \{0,1\}^N$ is then seen to be the lexicographically maximal sequence satisfying (\ref{eq:greedy exp for x}), which exactly corresponds to the first $N$ bits of the greedy $\beta$-expansion of $R_i + s_i + \varepsilon$, i.e.
\begin{equation}\label{eq:def x n i x n i + 1}
	x := d_{\beta_{i+1}}^N(R_i + s_i + \varepsilon_i).
\end{equation}
Therefore, $\xf$ is iteratively defined as 
\begin{equation}
	\xf = \coprod_{i=0}^\infty d_{\beta_{i+1}}^N\left(R_i + s_i+\varepsilon_i\right).
\end{equation}
This does not match (\ref{eq:inferring chunks of greedy expansion from R i E beta}), and
hence, as opposed to the simpler the case $\beta \in \Q$, we cannot conclude on the relation between $R_{i+1}$ and $R_i$. Instead, we establish the relation as follows. For $i \in \N_0$, we have
\allowdisplaybreaks
\begin{align}
	R_{i+1} &\overset{(\ref{eq:R i definition beta r})}{=}\beta_{i+2}^{N(i+1)}\left(\delta_2(\gf_2(s)_{1:\Sigma_{i+1}}) - \delta_{\beta_{i+1}}(\xf_{1:N(i+1)})\right)\\
	& \osetlem{\ref{lem:factorisation of beta expansion through concatenation two}}{=} \beta_{i+2}^{N(i+1)}\left(\delta_2(\gf_2(s)_{1:\Sigma_{i}}) + \frac{1}{2^{\Sigma_i}}\delta_2(\gf_2(s)_{\Sigma_i+1:\Sigma_{i+1}}) - \delta_{\beta_{i+1}}(\xf_{1:Ni}) \right.\nonumber\\
	&  \hspace{2.3cm}\left. - \frac{1}{\beta_{i+1}^{Ni}}\delta_{\beta_{i+1}}(\xf_{Ni+1:N(i+1)}) \right)\\
	& = \beta_{i+2}^{N(i+1)}\left(\delta_2(\gf_2(s)_{1:\Sigma_{i}}) + \frac{1}{2^{\Sigma_i}}\delta_2(\gf_2(s)_{\Sigma_i+1:\Sigma_{i+1}}) - \delta_{\beta_{i+1}}(\xf_{1:Ni}) \right.\nonumber\\
	&  \hspace{2.3cm} \left. - \delta_{\beta_i}(\xf_{1:Ni}) + \delta_{\beta_i}(\xf_{1:Ni}) - \frac{1}{\beta_{i+1}^{Ni}}\delta_{\beta_{i+1}}(\xf_{Ni+1:N(i+1)}) \right)\\
	& = \frac{\beta_{i+2}^{N(i+1)}}{\beta_{i+1}^{Ni}}\left(\beta_{i+1}^{Ni}\left(\delta_2(\gf_2(s)_{1:\Sigma_{i}}) - \delta_{\beta_i}(\xf_{1:Ni})\right) + \frac{\beta_{i+1}^{Ni}}{2^{\Sigma_i}}\delta_2(\gf_2(s)_{\Sigma_i+1:\Sigma_{i+1}})  \right.\nonumber\\
	&  \hspace{2.3cm} \left. + \beta_{i+1}^{Ni}\left(\delta_{\beta_i}(\xf_{1:Ni}) - \delta_{\beta_{i+1}}(\xf_{1:Ni})\right)  - \delta_{\beta_{i+1}}(\xf_{Ni+1:N(i+1)}) \right)\\
	& \overset{\substack{(\ref{eq:R i definition beta r})\\ (\ref{eq:s i to generate beta expansion beta r})\\ (\ref{eq:epsilon i to generate beta expansion beta r})}}{=} \frac{\beta_{i+2}^{N(i+1)}}{\beta_{i+1}^{Ni}}\left(R_i + s_i + \varepsilon_i  - \delta_{\beta_{i+1}}\left(\xf_{Ni+1:N(i+1)}\right) \right)\\
	& \overset{(\ref{eq:def x n i x n i + 1})}{=} \frac{\beta_{i+2}^{N(i+1)}}{\beta_{i+1}^{Ni}}\left(R_i + s_i + \varepsilon_i  - \delta_{\beta_{i+1}}\left(d_{\beta_{i+1}}^N\left(R_i + s_i + \varepsilon_i\right)\right) \right)\\
	&\overset{(\ref{eq:n iterations of greedy map is same as reading the first n bits})}{=} \frac{\beta_{i+2}^{N(i+1)}}{\beta_{i+1}^{Ni}} \frac{1}{\beta_{i+1}^{N}} G^N_{\beta_{i+1}}\left( R_i + s_i + \varepsilon_i\right).
\end{align}
We conclude on the relation
\begin{equation}
	R_{i+1} = \left(\frac{\beta_{i+2}}{\beta_{i+1}}\right)^{N(i+1)} G^N_{\beta_{i+1}}\left( R_i + s_i + \varepsilon_i\right).
\end{equation}
Note that in the above equation, $\varepsilon_i$ depends on $\xf_{1:Ni}$. The dynamical system hence needs to record the successive values taken by $\xf_{1:Ni}$. For that purpose, we add an extra variable, denoted by $X_i$, such that $X_i = \xf_{1:Ni}$. The dynamical system will hence be defined by a system of equations instead of just one equation.
\begin{equation}\label{eq:R i X i + 1 definition of E N}
	\begin{array}{rr}
		R_{i+1} =& \left(\frac{\beta_{i+2}}{\beta_{i+1}}\right)^{N(i+1)} G^N_{\beta_{i+1}}\left( R_i + s_i + \varepsilon_N(X_i, \beta_i, \beta_{i+1} , i)\right),\\
	X_{i+1} =& X_i \, d_{\beta_{i+1}}^N\left(R_i +  s_i + \varepsilon_N(X_i, \beta_i, \beta_{i+1} , i)\right).
	\end{array}
	\tag{\text{$\mathcal E_N$}}
\end{equation}
where $\varepsilon_N(x, t_1, t_2 , i) := t_2^{Ni}\left(\delta_{t_1}(x) - \delta_{t_2}(x)\right)$, for all $x \in \{0,1\}^\ast$, $t_1,t_2 \in (1,2)$ and $i \in \N_0$.
Similarly as for $E_\beta^N$, when presented with inputs $\ssf = s_0s_1\ldots \in  I_\beta^\omega$ and $\betaf = \beta_0\beta_1\ldots \in  (1,2)^\omega$, we can see $\EE_N$ as a mapping $ \bigcup_{n \in \N_0} \left(I_\beta^n \times (1,2)^n\right) \to \{0,1\}^n$ by $\EE_N(\epsilon, \epsilon) := \epsilon$ and
\begin{equation}\label{eq:inferring prefix of chunks of greedy expansion from R i E beta i}
	\EE_N(s_0s_1\ldots s_{n-1}, \beta_0\beta_1\ldots \beta_n) := X_{n-1} = \coprod_{i=0}^{n-1} d_{\beta_{i+1}}^N\left(R_i +  s_i + \varepsilon_N(X_i, \beta_i, \beta_{i+1} , i)\right), 
\end{equation}
for all $n \in \N$, and as a mapping $I_\beta^\omega \times (1,2)^\omega \to \{0,1\}^\omega$ through
\begin{equation}\label{eq:inferring chunks of greedy expansion from R i E beta i}
	\EE_N(\ssf, \betaf) := \limi{n} X_n = \coprod_{i=0}^{\infty} d_{\beta_{i+1}}^N\left(R_i +  s_i + \varepsilon_N(X_i, \beta_i, \beta_{i+1} , i)\right).
\end{equation}
One must pay attention to the fact that if $R_i + s_i + \varepsilon_N(X_i, \beta_i, \beta_{i+1} , i) \notin I_{\beta_{i+1}}$, then (\ref{eq:R i X i + 1 definition of E N}) is not well-defined as $G_{\beta_{i+1}}$ is defined on $I_{\beta_{i+1}}$ only. Using the interpretation in terms of a dynamical system, we establish conditions on $N,\Sigma_i, \Lambda_i$ so that, given $\beta$, (\ref{eq:R i X i + 1 definition of E N}) is well-defined, then prove that the above procedure produces a valid $\beta$-expansion, and finally study the algorithmic complexity of the delivered $\beta$-expansion.

\subsection{The dynamical system is well-defined}

For (\ref{eq:R i X i + 1 definition of E N}) to be well-defined, we must have $R_i + s_i + \varepsilon_N(X_i, \beta_i, \beta_{i+1} , i) \in I_{\beta_{i+1}}$ for all $i \in \N_0$. Note that $\beta_i \leq \beta$ implies $I_{\beta} \subset I_{\beta_{i+1}}$, for all $i \in \N_0$. Therefore, in this case, this is sufficient to impose $R_i + s_i + \varepsilon_N(X_i, \beta_i, \beta_{i+1} , i) \in I_{\beta}$ for all $i \in \N_0$. We define
\begin{equation}
	\dom \EE_N(\beta) := \left\{ \ssf \in I_\beta^\omega, \betaf \in (1,\beta]^\omega : R_i + s_i + \varepsilon_N(X_i, \beta_i, \beta_{i+1} , i) \in I_{\beta}, \ \forall i \in \N_0  \right\}.
\end{equation}
The structure of the proof is almost identical to the proof of the well-definition of (\ref{eq:greedy map chunks as dynamical system}) for the simpler case $\beta \in \Q$. Concretely, we set $K_\beta := \frac{1}{3}J_\beta$, where $J_\beta \neq \varnothing$ is defined by (\ref{eq:definition J beta}), for $\beta \in (1,2)$, such that $I'_\beta := [0,1] + 3K_\beta = [0, \frac{\beta/2}{\beta-1}] \subset [0, \frac{1}{\beta-1}] = I_\beta$. We will prove, under certain conditions on $N$, $\Sigma_i$ and $\Lambda_i$, that $R_i \in [0,1] + K_\beta$, $s_i \in K_{\beta}$ and $\varepsilon_N(X_i, \beta_i, \beta_{i+1} , i) \in K_{\beta}$, resulting in
\begin{equation}
	\underset{\in [0,1]+K_\beta}{\underbrace{R_i}} + \underset{\in K_\beta}{\underbrace{s_i}} + \underset{\in K_\beta}{\underbrace{\varepsilon_N(X_i, \beta_i, \beta_{i+1} , i)}} \in [0,1] + 3K_\beta \subset I_\beta.
\end{equation}
We start by proving that under some conditions on $N$, $\ssf$ and $\betaf$, $\EE_N$ is well-defined. For $\beta \in (1,2)$, we define $\lowbeta := 1 + \frac{\beta-1}{10} > 1$ and introduce the set $A_\beta = \left[\lowbeta, \beta \right] \subset (1,\beta)$, displayed on Figure \ref{fig:set A beta}, that will turn out very useful in the upcoming discussion.
\begin{figure}[H]
	\centering
	\begin{tikzpicture}[scale = 5]
		\draw (0,0) -- (2,0);
		\draw (0,-0.02) -- (0,0.02);
		\draw (1,-0.02) -- (1,0.02);
		\draw (2,-0.02) -- (2,0.02);
		\draw[blue] (1.06,-0.02) -- (1.06,0.02);
		\draw[blue] (1.6,-0.02) -- (1.6,0.02);
		\draw[blue, thick] (1.06,0) -- (1.6,0);
		\node[below] at (0,-0.015){$0$};
		\node[below] at (1,-0.015){$1$};
		\node[below] at (2,-0.015){$2$};
		\node[above] at (1.12,-0.012){\tiny $1 + \frac{\beta-1}{10}$};
		\node[above] at (1.6,-0.008){\tiny $\beta$};
		\node[below] at (1.3,-0.008){\footnotesize $A_\beta$};
	\end{tikzpicture}
	\caption{Example of $A_\beta$ for $\beta = 1.6$.}
	\label{fig:set A beta}
\end{figure}
We first make explicit a subset of $I_\beta^\omega \times (1, 2)^\omega$ taht is included in $\dom \EE_N(\beta)$.
\begin{lemma}\label{eq:J beta is in dom E beta N beta r}
	Let $\beta \in (1,2)$. If 
	\begin{equation}\label{eq:condition on N}
		\tag{$\mathcal N$}
		N \geq \log_{\lowbeta}(2) + \log_{\lowbeta}\left( \frac{2- {\lowbeta}}{2 - \beta} \right),
	\end{equation}
	then
	\begin{equation}
		K_\beta^\omega \times \mathcal U_N(\beta) \subset \dom \EE_N(\beta),
	\end{equation}
	where $\mathcal U_N(\beta) := A_\beta^\omega \cap \mathcal V_N(\beta) \cap \mathcal W_N(\beta)$, with
	\begin{equation}\label{eq:definition V beta N}
		\mathcal V_N(\beta) := \left\{ \betaf \in (1,2]^\omega : \varepsilon_N(x, \beta_i, \beta_{i+1} , i) \in K_{\beta}, \ \text{for all} \ x \in \{0,1\}^{Ni}, \ i \in \N_0 \right\},
	\end{equation}
	and
	\begin{equation}\label{eq:definition W beta N}
		\mathcal W_N(\beta) := \left\{ \betaf \in (1,2]^\omega : \left(\frac{\beta_{i+1}}{\beta_{i}}\right)^{Ni} \in [0,1] + K_{\beta}, \ i \in \N_0 \right\}.
	\end{equation}
\end{lemma}
\begin{proof} 
	Let $\beta \in (1,2)$, $N \geq \log_\beta(2)$, $\ssf = s_0s_1s_2\ldots \in K_\beta^\omega$ and $\betaf =\beta_0\beta_1\ldots \in \mathcal U_N(\beta)$. We prove by induction that
	\begin{equation}\label{eq:condition Ri si beta i}
		R_i + s_i + \varepsilon_N(X_i, \beta_i, \beta_{i+1} , i) \in I'_\beta = [0,1] + 3K_\beta \subset I_\beta, \ \text{for all}\ i \in \N_0.
	\end{equation}
	We start by showing that $R_0 + s_0 + \varepsilon_N(X_0, \beta_0, \beta_{1} , 0) \in I'_\beta$. Indeed, from the definition of $\EE_N$, $R_0 = 0$. Moreover, as $s_0 \in K_\beta$ and $\betaf \in \mathcal V_N(\beta)$, one has
	\begin{equation}
		R_0 + s_0 + \varepsilon_N(X_0, \beta_0, \beta_{0+1} , 0) = s_0 + \varepsilon_N(X_0, \beta_0, \beta_{1} , 0) \in K_\beta + K_\beta \subset I'_\beta.
	\end{equation}
	We continue by proving that if $S_{i-1} := R_{i-1} + s_{i-1} + \varepsilon_N(X_{i-1}, \beta_{i-1}, \beta_{i} , i-1)  \in I'_\beta$, then $S_i := R_{i} + s_i+ \varepsilon_N(X_i, \beta_i, \beta_{i+1} , i) \in I'_\beta$, for all $i \in \N$. Fix $i \in \N$ and suppose that $S_{i-1}\in I'_\beta$. It follows from (\ref{eq:R i X i + 1 definition of E N}) that
	\begin{equation}\label{eq:relation R i+1 R i s i beta i}
		R_i = \left(\frac{\beta_{i+1}}{\beta_{i}}\right)^{Ni} G_{\beta_{i+1}}^N\left(S_{i-1}\right).
	\end{equation}
	By Lemma \ref{ineq:greedy beta expansion}, 
	\begin{equation}
		G_{\beta_{i+1}}^n\left(S_{i-1}\right) \in [0,1]
	\end{equation}
	for all
	\begin{equation}\label{eq:condition on n for convergence of greedy expansion beta r}
		n \geq - \log_{\beta_{i+1}}\left( \frac{1 - ({\beta_{i+1}} - 1)S_{i-1}}{2- {\beta_{i+1}}} \right).
	\end{equation}
	From $S_{i-1} \in I'_\beta = \left[0, \frac{\beta/2}{\beta-1}\right]$, one calculates
	\begin{align}
		- \log_{\beta_{i+1}}\left( \frac{1 - ({\beta_{i+1}} - 1)S_{i-1}}{2- {\beta_{i+1}}} \right) &\leq - \log_{\beta_{i+1}}\left( \frac{1 - ({\beta_{i+1}} - 1)\frac{\beta/2}{\beta-1}}{2- {\beta_{i+1}}} \right)\overset{(a)}{\leq} - \log_{\beta_{i+1}}\left( \frac{1 - \beta/2}{2- {\beta_{i+1}}} \right)\nonumber\\
		& = \log_{\beta_{i+1}}(2) + \log_{\beta_{i+1}}\left( \frac{2- {\beta_{i+1}}}{2 - \beta} \right)\\
		&\overset{(b)}{\leq} \log_{\lowbeta}(2) + \log_{\lowbeta}\left( \frac{2- {\lowbeta}}{2 - \beta} \right)\\
		& \leq N,
	\end{align}
	where (a) and (b) use that $\beta_{i+1} \in A_\beta = [\lowbeta, \beta]$
	Then, $N$ satisfies (\ref{eq:condition on n for convergence of greedy expansion beta r}), so
	\allowdisplaybreaks
	\begin{align}\label{eq:R i is in 0 1 beta r}
		R_{i} &\overset{(\ref{eq:relation R i+1 R i s i beta i})}{=} \left(\frac{\beta_{i+1}}{\beta_{i}}\right)^{Ni} G^N_\beta(R_{i-1} + s_{i-1} + \varepsilon_N(X_{i-1}, \beta_{i-1}, \beta_{i} , i-1))\\
		& \osetlem{\ref{ineq:greedy beta expansion}}{\in} \left(\frac{\beta_{i+1}}{\beta_{i}}\right)^{Ni}[0,1] \overset{(a)}{=} [0,1] + K_\beta,
	\end{align}
	where (a) follows from $\betaf \in \mathcal W_N(\beta)$, which implies $\left(\frac{\beta_{i+1}}{\beta_{i}}\right)^{Ni} \in [0,1] + K_\beta$. Moreover, $s_i \in K_\beta$ and $\betaf \in \mathcal V_N(\beta)$, so we have $\varepsilon_N(X_i, \beta_i, \beta_{i+1} , i) \in K_\beta$. It follows that $R_i + s_i + \varepsilon_N(X_i, \beta_i, \beta_{i+1} , i) \in [0,1] + 3 K_\beta = I'_\beta$. By virtue of inductive reasonning, we conclude that $R_i + s_i + \varepsilon_N(X_i, \beta_i, \beta_{i+1} , i) \in I'_\beta$ for all $i \in \N_0$, so $(\ssf,\betaf) \in \dom \EE_N$.
\end{proof}
We now identify explicit subsets of $\mathcal V_N(\beta)$ and $\mathcal W_N(\beta)$. We start by making explicit a subset $\mathcal V'_N(\beta) \subset \mathcal V_N(\beta)$. For $\beta \in (1,2)$, define $C_\beta := \max K_\beta$.

\begin{definition} Let $\beta \in (1,2)$ and $N \in \N$. We say that $\betaf \in (1,\beta)^\omega$ converges $N$-fast towards $\beta$ if 
	\begin{equation}\label{eq:C N fast converging}
		|\beta - \beta_i| \leq \frac{\min\{1,C_\beta\}}{e\beta^{Ni}N^2i^2}, \ \text{for all} \ i \in \N.
	\end{equation}
We denote by $\mathcal Z_N(\beta)$ the set of nondecreasing sequences $\betaf \in (1,\beta)^\omega$ that converge $N$-fast towards $\beta$.
\end{definition}
\begin{lemma}\label{lem:beta C N converging implies bounded error}
	For $\beta \in (1,2)$ and $N \in \N$, $\mathcal Z_N(\beta)\subset \mathcal V_N(\beta) \cap W_N(\beta)$. 
\end{lemma}
\begin{proof} Let $\beta \in (1,2)$, $N \in \N$ and $\betaf \in \mathcal Z_N(\beta)$. We first show that $\mathcal Z_N(\beta) \subset \mathcal V_N(\beta)$, and then move to the proof of $\mathcal Z_N(\beta) \subset \mathcal W_N(\beta)$. Fix $i \in \N_0$ and $x \in \{0,1\}^{Ni}$. We will show that $\betaf \in \mathcal V_N(\beta)$ by proving that $\varepsilon_N(x, \beta_i, \beta_{i+1} , i) \in K_\beta =: [0, C_\beta]$. First note that $\varepsilon_N(\epsilon, \beta_0, \beta_{1} , 0) = 0 \in K_\beta$, so the case $i = 0$ is dealt with. Suppose in the following that $i \geq 1$. As $\betaf$ is nondecreasing, $\beta_i \leq \beta_{i+1}$ so $\delta_{\beta_i}(x) \geq \delta_{\beta_{i+1}}(x)$, and hence $0 \leq \beta_{i+1}^{Ni}\left(\delta_{\beta_i}(x) - \delta_{\beta_{i+1}}(x)\right) = \varepsilon_N(x, \beta_i, \beta_{i+1} , i)$. Let us now prove that $\varepsilon_N(x, \beta_i, \beta_{i+1} , i) \leq C_\beta$.
	\begin{align}
		\delta_{\beta_i}(x) - \delta_{\beta_{i+1}}(x) &= \sum_{j=1}^{Ni} x_j \beta_i^{-j} - \sum_{j=1}^{Ni} x_j \beta_{i+1}^{-j} = \sum_{j=1}^{Ni} x_j \left(\beta_i^{-j} -  \beta_{i+1}^{-j}\right)\\
		&\leq \sum_{j=1}^{Ni} x_j  \left(\beta_i^{-j} - \beta_{i+1}^{-j}\right) \overset{(a)}{\leq} \sum_{j=1}^{Ni}  \left(\beta_i^{-j} - \beta_{i+1}^{-j}\right), \label{eq:weak eq bounding delta beta variations}
	\end{align}
	where (a) is by $x_i \leq 1$.
	Note that for all $\beta \in [1,2)$, $\left|\frac{\partial \beta^{-j}}{\partial \beta}\right| = -j \beta^{-(j+1)} \geq -j$, for all $j \in \N$. Then, by using the mean value theorem, one gets $\beta_{i}^{-j} - \beta_{i+1}^{-j} \leq j (\beta_i - \beta_{i+1})$ for all $j \in \N$. By (\ref{eq:weak eq bounding delta beta variations}),
	\begin{align}
		\beta_{i+1}^{Ni}\left(\delta_{\beta_i}(x) - \delta_{\beta_{i+1}}(x)\right) &\leq \beta_{i+1}^{Ni} \sum_{j=1}^{Ni} j (\beta_i - \beta_{i+1}) = \beta_{i+1}^{Ni} \frac{Ni(Ni+1)}{2} (\beta_i - \beta_{i+1})\\
		&\leq \beta_{i+1}^{Ni} N^2i^2 (\beta_i - \beta_{i+1}) \overset{(a)}{\leq} \beta_{i+1}^{Ni} N^2i^2 (\beta - \beta_{i+1}) \overset{(\ref{eq:C N fast converging})}{\leq} C_\beta,
	\end{align}
	where (a) follows from $\betaf$ being nondecreasing. This finishes the proof that $\mathcal Z_N(\beta) \subset \mathcal V_N(\beta)$. We now prove that $\mathcal Z_N(\beta) \subset \mathcal W_N(\beta)$. Fix $i \in \N_0$ and $x \in \{0,1\}^{Ni}$. We will show that $\betaf \in \mathcal W_N(\beta)$ by proving that $\left(\frac{\beta_{i+1}}{\beta_{i}}\right)^{Ni} \in [0,1] + K_\beta =: [0, 1+ C_\beta]$. First note that $\left(\frac{\beta_{1}}{\beta_{0}}\right)^{0} = 1 \in [0, 1+ C_\beta]$, so the case $i = 0$ is dealt with. Suppose in the following that $i \geq 1$. (\ref{eq:C N fast converging}) implies that, in particular,
	\begin{equation}\label{eq:weak N fast convergence}
		|\beta - \beta_i| \leq \frac{\min\{1,C_\beta\}}{eNi}\cdot
	\end{equation}
	Then,
	\begin{align}
		\left(\frac{\beta_{i+1}}{\beta_{i}}\right)^{Ni} &= \left(1 + \frac{\beta_{i+1} - \beta_i }{\beta_{i}}\right)^{Ni} \overset{(a)}{<}  \left(1 + |\beta - \beta_i|\right)^{Ni}  \overset{(b)}{\leq} 1 + eNi |\beta - \beta_i| \overset{(\ref{eq:weak N fast convergence})}{\leq} 1 + C_\beta,
	\end{align}
	where (a) is by $\beta_i > 1$ and $\beta_{i+1} \leq \beta$, and (b) follows from the study of the sign of $(1+t)^n - (1 + ent)$ for $x \in [0,1]$ and $n \in \N$. This concludes the proof of $\mathcal Z_N(\beta) \subset \mathcal W_N(\beta)$.
\end{proof}
We now establish explicit conditions on $\Lambda_i$ so that $\betaf$ converges $N$-fast towards $\beta$.

\begin{lemma}\label{lem:condition on L for beta C N converging}
	Let $\beta \in (1,2)$, $N,L > 0$, and define $\Lambda_i := Ni + L$ for all $i \in \N_0$, and define $\betaf = \beta_0\beta_1\ldots$ by $\beta_i := 1 + \delta_2(\gf_2(\betaun)_{\Lambda_i})$ for all $i \in \N_0$. Suppose that 
	\begin{equation}\label{eq:condition on L}
		\tag{\text{$\mathcal L_1$}}
		L \geq 1 -\log_2\left(\min\{1,C_\beta\} (1-\log_2(\beta)^2)\right).
	\end{equation}
	Then, $\betaf$ is nondecreasing and converges $N$-fast towards $\beta$.
\end{lemma}
\begin{proof}
	Let $\beta \in (1,2)$, $N,L > 0$, and define $\Lambda_i := Ni + L$ for all $i \in \N_0$, and define $\betaf = \beta_0\beta_1\ldots$ by $\beta_i := 1 + \delta_2(\gf_2(\betaun)_{\Lambda_i})$ for all $i \in \N_0$. $\betaf$ is obviously nondecreasing. Moreover,
	\begin{equation}
		|\beta - \beta_i| = |\beta - 1 - \delta_2(\gf_2(\betaun)_{\Lambda_{i}})| \leq \frac{1}{2^{\Lambda_{i}}}.
	\end{equation}
	We will prove that
	\begin{equation} \label{eq:condition lambda i for C N fast}
		\frac{1}{2^{\Lambda_{i}}} \leq \frac{\min\{1,C_\beta\}}{e\beta^{Ni}N^2i^2}, \ \text{for all} \ i \in \N_0,
	\end{equation}
	which will alow us to conclude that $\betaf$ converges $N$-fast. 
	\begin{align}
		\frac{1}{2^{\Lambda_{i}}} = \frac{1}{2^{Ni+L}} &\overset{(\ref{eq:condition on L})}{\leq} \frac{1}{2^{Ni} 2^{1 - \log_2(\min\{1,C_\beta\}(1-\log_2(\beta)^2))}}\\ 
		&= \frac{\min\{1,C_\beta\} (1-\log_2(\beta))^2}{2\times2^{Ni}eN}\\
		&= \frac{eN^2(1-\log_2(\beta))^2\beta^{Ni}i^2}{2\times 2^{Ni}}\frac{\min\{1,C_\beta\}}{\beta^{Ni}eN^2i^2}\cdot
	\end{align}
	The function $f : t \mapsto \frac{\beta^{Nt}t^2}{2^{Nt}}$ attains its maximum in $t_0 = \frac{2\ln(2)}{N(1-\log_2(\beta))}$, of a value of
	\begin{equation}
		f(t_0) = \frac{4\ln(2)^2}{2^{2\ln(2)} N^2 (1-\log_2(\beta))^2}.
	\end{equation}
	Therefore,
	\begin{align}
			\frac{1}{2^{\Lambda_{i}}} &\leq \frac{e \times 4\ln(2)^2}{2 \times 2^{2\ln(2)}}\frac{\min\{1,C_\beta\}}{\beta^{Ni}eN^2i^2}\cdot
	\end{align}
	The proof is concluded by remarking that $\frac{e \times 4\ln(2)^2}{2 \times 2^{2\ln(2)}} \simeq 0.999 \leq 1$.
	% (\ref{eq:condition lambda i for C N fast}) is equivalent to
	% \begin{equation}
	% 	Ni + L = \Lambda_i \geq Ni \log_2(\beta) + 2 \log_2(i) + \log_2(eN^2) - \log_2(\min\{1,C_\beta\}), \ \text{for all} \ i \in \N_0,
	% \end{equation}
	% which holds if
	% \begin{equation}
	% 	L \geq \max_{i\in\N_0} \left\{ Ni(\log_2(\beta) - 1) + 2\log_2(i)\right\} - \log_2(C_\beta) .
	% \end{equation}
	% The function $f : t \mapsto Nt (\log_2(\beta) - 1) + 2 \log_2(t)$ 
\end{proof}
\noindent We further establish explicit conditions on $\Lambda_i$ so that $\betaf \in A_\beta^\omega$.
\begin{lemma}\label{lem:condition on L for beta in A beta}
	Let $\beta \in (1,2)$, $N,L > 0$, and define $\Lambda_i := Ni + L$ for all $i \in \N_0$, and define $\betaf = \beta_0\beta_1\ldots$ by $\beta_i := 1 + \delta_2(\gf_2(\betaun)_{\Lambda_i})$ for all $i \in \N_0$. Suppose that 
	\begin{equation}\label{eq:condition on L 2}
		\tag{\text{$\mathcal L_2$}}
		L \geq - \log_2\left(\frac{9}{10}(\beta-1)\right).
	\end{equation}
	Then, $\betaf \in A_\beta^\omega$.
\end{lemma}
\begin{proof}
	Let $\beta \in (1,2)$, $N,L > 0$, and define $\Lambda_i := Ni + L$ for all $i \in \N_0$, and define $\betaf = \beta_0\beta_1\ldots$ by $\beta_i := 1 + \delta_2(\gf_2(\betaun)_{\Lambda_i})$ for all $i \in \N_0$.  Suppose that (\ref{eq:condition on L 2}) holds. Note that $\betaf$ is nondecreasing and satisfies $\beta_i \leq \beta$ for all $i \in \N_0$. Then, we only have to prove that $\beta_0 \geq \lowbeta$. Then,
	\begin{align}
		\beta_0 = 1 + \delta_2(\gf_2(\betaun)_{\Lambda_0}) &\geq 1 + \beta - 1 - 2^{-L} \geq \beta - 2^{\log_2\left(\frac{9}{10}(\beta-1)\right)} = 1 + \frac{\beta-1}{10} = \lowbeta . 
	\end{align}
\end{proof}

We finally conclude with concrete conditions on both the $\Sigma_i$ and $\Lambda_i$ so that the dynamical system is well-defined.

% \begin{theorem}
% 	Let $\beta \in (1,2)$, $N,L,C > 0$, and define $\Lambda_i := Ni + L$ for all $i \in \N_0$, and define $\betaf = \beta_0\beta_1\ldots$ by $\beta_i := 1 + \delta_2(\gf_2(\betaun)_{\Lambda_i})$ for all $i \in \N_0$. Suppose that (\ref{eq:condition on L}) holds and that $C \in J_\beta$. Then, 
% 	\begin{equation}
% 		\beta_{i+1}^{Ni}\left(\delta_{\beta_i}(\xf_{1:i}) - \delta_{\beta_{i+1}}(\xf_{1:i})\right) \in J_{\beta_{i+1}}, \ \text{for all} \ i \in \N_0.
% 	\end{equation}
% \end{theorem}
% \begin{proof}
% 	Let $\beta \in (1,2)$, $N,L,C > 0$, and define $\Lambda_i := Ni + L$ for all $i \in \N_0$, and define $\betaf = \beta_0\beta_1\ldots$ by $\beta_i := 1 + \delta_2(\gf_2(\betaun)_{\Lambda_i})$ for all $i \in \N_0$. Suppose that (\ref{eq:condition on L}) holds and that $C \in J_\beta$. Since (\ref{eq:condition on L}), by Lemma \ref{lem:condition on L for beta C N converging} $\betaf$ is nondecreasing and converges $(C,N)$-fast. By Lemma \ref{lem:beta C N converging implies bounded error}, we deduce that 
% 	\begin{equation}
% 		0 \leq \beta_{i+1}^{Ni}\left(\delta_{\beta_i}(\xf_{1:i}) - \delta_{\beta_{i+1}}(\xf_{1:i})\right)\leq C.
% 	\end{equation}
% 	Since $C \in J_\beta$, it follows that $\beta_{i+1}^{Ni}\left(\delta_{\beta_i}(\xf_{1:i}) - \delta_{\beta_{i+1}}(\xf_{1:i})\right) \in J_\beta$. We conclude the proof by noticing that $J_\beta = \cap_{i \in \N_0} J_{\beta_{i+1}}$.
% \end{proof}

% As a conclusion of this section, we derive the following statement.
\begin{theorem}\label{thm:final conditions P1 P2}
	Let $s \in [0,1]$, $\beta \in (1,2)$, and
	\begin{equation}\label{eq:final value parameters beta r 1}\tag{\text{$\mathcal{P}_1$}}
		\begin{array}{rl}
			N &:= \left\lceil\log_{\lowbeta}(2) + \log_{\lowbeta}\left( \frac{2- {\lowbeta}}{2 - \beta} \right)\right\rceil,\\
			L &:= \left\lceil 1-\log_2\left(\frac{9}{10}\min\{1,C_\beta\}(\beta-1) (1-\log_2(\beta)^2)\right)\right\rceil\\
			\Lambda_i &:= Ni + L, \ \text{for all} \ i \in \N_0. 
		\end{array}
	\end{equation}
	We then define $\betaf = \beta_0\beta_1\ldots \in A_\beta$ according to (\ref{eq:beta i to generate beta expansion beta r}), $\Sigma_0 =0$, and
	\begin{equation}\label{eq:final value parameters beta r 2}\tag{\text{$\mathcal{P}_2$}}
		\Sigma_i := \lceil Ni\log_2(\beta_{i+1})\rceil - \left\lfloor \log_2\left(C_\beta\right) \right\rfloor, \ \text{for all} \ i \in \N.
	\end{equation}
	Finally, let $\ssf = s_0s_1\ldots$ be defined according to (\ref{eq:s i to generate beta expansion beta r}). Then, $(\ssf,\betaf) \in \dom \EE_N$.
\end{theorem}
\begin{proof}
	Let $\beta,s,N,L,\Lambda_i,\betaf, \Sigma_i$ and $\ssf$ be defined as above. We prove that $\betaf \in \mathcal U_N(\beta)$ and $\ssf \in K_\beta^\omega$, which allows to conclude by Lemma \ref{eq:J beta is in dom E beta N beta r}. We first show that $\betaf \in \mathcal U_N(\beta)$. Note that $L$ satisfies both (\ref{eq:condition on L}) and (\ref{eq:condition on L 2}), hence by Lemmata \ref{lem:condition on L for beta C N converging} and \ref{lem:condition on L for beta in A beta}, $\betaf \in A^\omega_\beta \cap \mathcal Z_N(\beta)$, which by Lemma \ref{lem:beta C N converging implies bounded error} implies $\betaf \subset A^\omega_\beta \cap \mathcal V_N(\beta)\cap \mathcal W_N(\beta)= \mathcal U_N(\beta)$. Now, show that $\ssf \in K_\beta^\omega$. Let $i \in \N_0$.
	\begin{align}
		s_i = \frac{\beta_{i+1}^{Ni}}{2^{\Sigma_i}} \delta_2(\gf_2(s)_{\Sigma_i+1:\Sigma_{i+1}}) \overset{(\ref{eq:final value parameters beta r 2})}{\leq} C_\beta \delta_2(\gf_2(s)_{\Sigma_i+1:\Sigma_{i+1}}) \overset{(a)}{\leq} C_\beta,
	\end{align}
	where (a) follows from $\delta_2(\cdot) \leq 1$. Hence, $s_i \in [0,C_\beta] = K_\beta$. Lemma \ref{eq:J beta is in dom E beta N beta r} yields $(\ssf, \betaf) \in K_\beta \times \UU_N(\beta) \subset \dom \EE_N(\beta)$.
\end{proof}

We now formally establish that the dynamical system generates a $\beta$-expansion.

\subsection{\texorpdfstring{The dynamical system produces a valid $\beta$-expansion}{}}


In this section, we show that for $\beta \in (1,2)$ and $s \in [0,1]$, if we define $N$, $L$ and $(\Lambda_i)_{i \in \N_0}$ according to (\ref{eq:final value parameters beta r 1}), $\betaf = \beta_0\beta_1\ldots$ according to (\ref{eq:beta i to generate beta expansion beta r}), $(\Sigma_i)_{i \in \N_0}$  and (\ref{eq:final value parameters beta r 2}),  and $\ssf = s_0s_1\ldots $ according to (\ref{eq:s i to generate beta expansion beta r}), then $\EE_N(\ssf, \betaf)$ is a $\beta$-expansion of $s$.

\begin{theorem}\label{thm:convergence rational dynamical system to beta expansion beta r}
	Let $\beta \in (1,2)$ and $s \in [0,1]$. Define $N$, $L$ and $(\Lambda_i)_{i \in \N_0}$ according to (\ref{eq:final value parameters beta r 1}), $\betaf = \beta_0\beta_1\ldots$ according to (\ref{eq:beta i to generate beta expansion beta r}), $(\Sigma_i)_{i \in \N_0}$ according to (\ref{eq:final value parameters beta r 2}),  and $\ssf = s_0s_1\ldots $ according to (\ref{eq:s i to generate beta expansion beta r}). Then,
	\begin{equation}
		\EE_N(\ssf,\betaf) \in \Sigma_\beta(s).
	\end{equation}
\end{theorem}
\begin{proof}
	Let $\beta \in (1,2)$ and $s \in [0,1]$. Define $N$, $L$ and $(\Lambda_i)_{i \in \N_0}$ according to (\ref{eq:final value parameters beta r 1}), $\betaf = \beta_0\beta_1\ldots$ according to (\ref{eq:beta i to generate beta expansion beta r}), $(\Sigma_i)_{i \in \N_0}$ according to (\ref{eq:final value parameters beta r 2}),  and $\ssf = s_0s_1\ldots $ according to (\ref{eq:s i to generate beta expansion beta r}). By Theorem \ref{thm:convergence rational dynamical system to beta expansion beta r}, $\EE_N(\ssf,\betaf)$ is well-defined. To prove that $\EE_N(\ssf,\betaf) \in \Sigma_\beta(s)$, by (\ref{eq:definition beta expansion}) we need to show that $\delta_\beta(\EE_N(\ssf,\betaf)) = s$. We let 
	\begin{align}
		x^{(n)} &:= \EE_N(s_0s_1\ldots s_n,\beta_0\beta_1\ldots \beta_n), \ \text{for} \ n \in \N,\\
		\xf &= \EE_N(\ssf,\betaf),
	\end{align}
	and we successively show 
	\begin{equation}\label{eq:prefixes of d theta converge to d theta beta r}
		\limi{n} \delta_\beta\left(x^{(n)}\right) = \delta_\beta(\xf),
	\end{equation}
	\begin{equation}\label{eq:delta beta and delta beta n converge}
		\limi{n} \delta_\beta\left(x^{(n)}\right) = \limi{n} \delta_{\beta_n}\left(x^{(n)}\right),
	\end{equation}
	and
	\begin{equation}\label{eq:prefixes of d theta converge to s beta r}
		\limi{n} \delta_{\beta_n}\left(x^{(n)}\right) = s,
	\end{equation}
	which combination leads to $\delta_\beta(\xf) = s$. We start by proving (\ref{eq:prefixes of d theta converge to d theta}). Let $n \in \N$. Note that
	\begin{align}
		\xf_{1:Nn} &\overset{(\ref{eq:inferring chunks of greedy expansion from R i E beta i})}{=} \left[\coprod_{i=0}^{\infty} d_{\beta_{i+1}}^N\left(R_i +  s_i + \varepsilon_N(X_i, \beta_i, \beta_{i+1} , i)\right)\right]_{1:Nn}\\
		&= \coprod_{i=0}^{n-1} d_{\beta_{i+1}}^N\left(R_i +  s_i + \varepsilon_N(X_i, \beta_i, \beta_{i+1} , i)\right)\\
		&\overset{(\ref{eq:inferring prefix of chunks of greedy expansion from R i E beta i})}{=} \delta_\beta\left(x^{(n)}\right).\label{eq:d theta prefix is a prefix of d theta beta r}
	\end{align}
	It follows that
	\begin{align}
		\delta_\beta(\xf) &= \delta_\beta(\xf_{1:Nn} \xf_{Nn+1:\infty})\\
		&\osetlem{\ref{lem:factorisation of beta expansion through concatenation two}}{=} \delta_\beta(\xf_{1:Nn}) + \beta^{-Nn} \delta_\beta(\xf_{Nn+1:\infty})\\
		&\overset{(a)}{\in} \delta_\beta(\xf_{1:Nn}) + \beta^{-Nn} I_\beta\\
		&\overset{(\ref{eq:d theta prefix is a prefix of d theta beta r})}{=} \delta_\beta\left(x^{(n)}\right)  + \beta^{-Nn}I_\beta,
	\end{align}
	where (a) is by $\delta_\beta(\yf) \in I_\beta$, for all $\yf \in \{0,1\}^\omega$. Hence, 
	\begin{equation}
		0 \leq \limi{n} \left(\delta_\beta(\xf) - \delta_\beta\left(x^{(n)}\right)\right) \leq \limi{n} \frac{\beta^{-Nn}}{\beta-1} = 0,
	\end{equation}
	yielding (\ref{eq:prefixes of d theta converge to d theta beta r}). We now move to the proof of (\ref{eq:delta beta and delta beta n converge}). Let $n \in \N$.
	\begin{align}
		\left|\delta_{\beta}\left(x^{(n)}\right) - \delta_{\beta_n}\left(x^{(n)}\right)\right| & \leq \sum_{i=1}^{Nn} x^{(n)}_i |\beta^{-i} - \beta_n^{-i}|\\
		& \overset{(a)}{\leq} \sum_{i=1}^{Nn} x^{(n)}_i i |\beta - \beta_n|\\
		& \leq \frac{Nn(Nn+1)}{2} |\beta - \beta_n|\\
		& \leq N^2n^2 |\beta - \beta_n|\\
		& \overset{(b)}{\leq} \frac{C_\beta}{\beta_n^{Nn}} \underset{n \to \infty}{\to} 0,
	\end{align}
	where (a) follows froms the arguments exposed just after (\ref{eq:weak eq bounding delta beta variations}) and (b) is a consequence of $\betaf \in \mathcal V'_N(\beta)$.
	
	We finally prove (\ref{eq:prefixes of d theta converge to s beta r}). Fix $n \in \N$ The dynamical system has been defined so that (\ref{eq:s i to generate beta expansion beta r}) holds, then
	\begin{equation}
		R_i \overset{(\ref{eq:R i definition beta r})}{=} \beta_{i+1}^{Ni}\left(\delta_2(\gf_2(s)_{1:\Sigma_{i}}) - \delta_{\beta_{i}}(\xf_{1:Ni})\right).
	\end{equation}
	As (\ref{eq:final value parameters beta r 1}) and (\ref{eq:final value parameters beta r 2}) hold, then by Theorem (\ref{thm:final conditions P1 P2}) $(\ssf, \betaf) \in \dom \EE_N$, i.e.,
	\begin{equation}R_n + s_n + \varepsilon_N(X_n, \beta_{n+1}, \beta_n, n) \in I_\beta.\end{equation}
	In particular, $0 \leq R_n \leq \frac{1}{\beta-1}$.
	\begin{equation}
		0 \leq \frac{R_i}{\beta^{Ni}} \overset{(\ref{eq:R i definition})}{=} \delta_2(\gf_2(s)_{\Sigma_i:\Sigma_{i+1}}) - \delta_\beta(x^{(n)})  \overset{(\ref{eq:R i definition})}{=} \frac{R_i}{\beta_n^{Ni}} \leq \frac{1}{(\beta-1) \beta_n^{Ni}} \overset{(a)}{\leq} \frac{1}{(\beta-1) \lowbeta^{Ni}}\underset{n \to \infty}{\to} 0,
	\end{equation}
	where (a) follows from $\betaf \in A_\beta^\omega$. Therefore,
	\begin{equation}
		\limi{n} \delta_{\beta_n}(x^{(n)}) = \limi{n} \delta_2(\gf_2(s)_{\Sigma_i:\Sigma_{i+1}}) = s.
	\end{equation}
	This ends the proof.
	% \allowdisplaybreaks
	% \begin{align}
		
	% 	\frac{R_n}{\beta_{n+1}} = 
	% \end{align}
	% \begin{align}
	% 	\delta_\beta(\EE_N(s_0s_1\ldots s_n,\beta_0\beta_1\ldots \beta_n)) &\overset{(\ref{eq:inferring prefix of chunks of greedy expansion from R i E beta i})}{=} \delta_\beta\left(\coprod_{i=0}^{n-1} d_{\beta_{i+1}}^N\left(R_i +  s_i + \varepsilon_N(X_i, \beta_i, \beta_{i+1} , i)\right)\right)\\
	% 	&\osetlem{\ref{lem:factorisation of beta expansion through concatenation}}{=}\sum_{i=0}^{n-1} \beta^{-Ni}\delta_\beta\left( d_{\beta_{i+1}}^N\left(R_i + s_i + \varepsilon_N(X_i, \beta_i, \beta_{i+1} , i)\right)\right)\\
	% 	&\overset{(\ref{eq:definition multiple digit map})}{=}\sum_{i=0}^{n-1} \beta^{-Ni}\delta_\beta\left(\gf_\beta\left(R_i + \frac{\beta^{Ni}}{2^{\Sigma_i}} s_i\right)_{1:N}\right)\\
	% 	&\overset{(\ref{eq:n iterations of greedy map is same as reading the first n bits})}{=} \sum_{i=0}^{n-1} \beta^{-Ni} \left( R_i + \frac{\beta^{Ni}}{2^{\Sigma_i}} s_i - \frac{1}{\beta^N}G^N_\beta\left(R_i + \frac{\beta^{Ni}}{2^{\Sigma_i}} s_i\right) \right)\\
	% 	&\overset{(\ref{eq:greedy map chunks as dynamical system})}{=} \sum_{i=0}^{n-1} \beta^{-Ni} \left(R_i + \frac{\beta^{Ni}}{2^{\Sigma_i}}s_i - \frac{1}{\beta^N}R_{i+1}\right)\\
	% 	&\overset{(\ref{eq:s i to generate beta expansion})}{=} \sum_{i=0}^{n-1} \beta^{-Ni} \left(R_i + \frac{\beta^{Ni}}{2^{\Sigma_i}}\delta_2(\gf_2(s)_{\Sigma_i+1:\Sigma_{i+1}}) - \frac{1}{\beta^N}R_{i+1}\right)\\
	% 	&= \sum_{i=0}^{n-1} \beta^{-Ni} R_i + \sum_{i=0}^{n-1} \frac{1}{2^{\Sigma_i}}\delta_2(\gf_2(s)_{\Sigma_i+1:\Sigma_{i+1}}) - \sum_{i=0}^{n-1} \beta^{-N(i+1)} R_{i+1}\\
	% 	&= R_0 - \beta^{Nn}R_n + \sum_{i=0}^{n-1} \frac{1}{2^{\Sigma_i}}\delta_2(\gf_2(s)_{\Sigma_i+1:\Sigma_{i+1}})\\
	% 	&\osetlem{\ref{lem:factorisation of beta expansion through concatenation}}{=} R_0 - \beta^{Nn}R_n + \delta_2(\gf_2(s)_{1:\Sigma_n}). \label{eq:intermediate equation delta beta x n}
	% \end{align}
	% From the definition of $E_\beta^N$, $R_0 = 0$. Moreover, as $(N, (\Sigma_i)_{i \in \N})$ satisfy (\ref{eq:final value parameters}), by (\ref{eq:R i is in 0 1}) one has $R_n \in [0,1]$. Therefore,
	% \begin{equation}
	% 	0 \geq \limi{n} \delta_\beta(E_\beta^N(s_0s_1 \ldots s_{n-1})) - \delta_2(\gf_2(s)_{1:\Sigma_n}) \geq \limi{n} -\beta^{-Nn} = 0.
	% \end{equation}
	% As $\limi{n} \Sigma_n = \infty$, we then have
	% \begin{equation}
	% 	\limi{n} \delta_\beta(E_\beta^N(s_0s_1 \ldots s_{n-1})) = \limi{n}\delta_2(\gf_2(s)_{1:\Sigma_n}) = s, 
	% \end{equation}
	% thereby finishing the proof.
\end{proof}

% \begin{theorem}\label{thm:generalized convergence rational dynamical system to beta expansion}
% 	Let $\beta \in (1,2)$, and suppose that $\DD_\beta$ satisfies $(\mathbf{H}')$. Then,
% 	\begin{equation}
% 		\xf_\beta(s) \in \Sigma_\beta(s), \ \forall s \in [0,1].
% 	\end{equation}
% \end{theorem}
% \begin{proof}
% 	Let $\beta \in (1,2)$, suppose that $\DD_\beta$ satisfies $(\mathbf{H}')$, and let $s \in [0,1]$. To prove that $\xf_\beta(s) \in \Sigma_\beta(s)$, we need to show that $\delta_\beta(\xf_\beta(s)) = s$. Consider the sequence $(X_n)_{n \in \N}$ generated by $\DD_\beta(s)$. $\limi{n} X_n = \xf_\beta(s)$, where the limit is defined according to distance $d$ introduced in Definition \ref{def:metric on discrete sequences}. As by Lemma \ref{lem:delta beta is continuous}, $\delta_\beta$ is continuous as a mapping from $\left(\{0,1\}^\omega, d\right)$ to $(\R,|\cdot|)$, $\left(\delta_\beta(X_n)\right)_{\N_2}$ therefore converges to $\delta_\beta(\xf_\beta(s))$. Further, fix $n \in \N$ and note that
% 	\begin{align}
% 		|\delta_\beta(X_n) - \delta_{q_{n-1}}(X_n)| &= \left|\sum_{i=1}^{|X_n|} [X_n]_i \left(\beta^{-i} - q_{n-1}^{-i}\right)\right| \leq \left|\sum_{i=1}^{|X_n|} \left(\beta^{-i} - q_{n-1}^{-i}\right)\right| \\
% 		&\leq \left|\sum_{i=1}^{\infty} \left(\beta^{-i} - q_{n-1}^{-i}\right)\right| \leq \frac{1}{\beta-1} - \frac{1}{q_{n-1}-1} \underset{n \to \infty}{\to} 0.
% 	\end{align}
% 	Then, $\limi{n}\delta_\beta(X_n) = \limi{n} \delta_{q_{n-1}}(X_n)$. We now show that  $\left(\delta_{q_{n-1}}(X_n)\right)_{\N_2}$ converges to $s$, which allows us to conclude that $ s = \limi{n}\delta_\beta(X_n) = \delta_\beta(\xf_\beta(s))$, and hence that $\delta_\beta(\xf_\beta(s)) \in \Sigma_\beta(s)$. Fix $n \in \N$. First note that $X_{n+1} = X_n  \gf_{q_{n}}(R'_{n})_{1: N}$. By repeating this argument, one has 
% 	\begin{equation}\label{eq:generalized B_n as concatenation of greedy Z i}
% 		X_{n+1} = X_0  \coprod_{i=0}^{n} \gf_{q_i}(R'_i)_{1: N} = \coprod_{i=0}^{n} \gf_{q_i}(R'_i)_{1: N}.
% 	\end{equation}
% 	Further, for $i \in \{1,\ldots,n\}$,
% 	\begin{align}
% 		R_{i+1} - q_i^N R_i &= q_i^N(1-\xi_i^N)R_i + \xi_i^N G_{q_i}^N(R'_i) - q_i^N R_i\\
% 		&\overset{(\ref{eq:n iterations of greedy map is same as reading the first n bits})}{=}q_i^N(1-\xi_i^N)R_i + \xi_i^N q_i^N\left(R'_i - \delta_{q_i}(\gf_{q_i}(R'_i))_{1: N}\right)- q_i^N R_i\\
% 		&=  \xi_i^N q_i^N\left(\frac{q_i^{Ni}}{2^{\Sigma_i}} S'_{i} + q_{i}^{Ni} \left(\delta_{q_{i-1}}(X_i) - \delta_{q_i}(X_i)\right) - \delta_{q_i}(\gf_{q_i}(R'_i))_{1: N}\right)\\
% 		&=  \chi_i^N \left(\frac{1}{2^{\Sigma_i}} S'_{i} + \delta_{q_{i-1}}(X_i) - \delta_{q_i}(X_i) - q_i^{-Ni}\delta_{q_i}(\gf_{q_i}(R'_i))_{1: N}\right)\\
% 		&\overset{(\ref{eq:generalized B_n as concatenation of greedy Z i})}{=}  \chi_i^N \left(\frac{1}{2^{\Sigma_i}} S'_{i} + \delta_{q_{i-1}}(X_i) - \delta_{q_i}\left(\coprod_{j=0}^{i-1} \gf_{q_j}(R'_j)_{1: N}\right)  - q_i^{-Ni}\delta_{q_i}(\gf_{q_i}(R'_i))_{1: N}\right)\nonumber \\
% 		&\overset{\ref{lem:factorisation of beta expansion through concatenation}}{=}  \chi_i^N \left(\frac{1}{2^{\Sigma_i}} S'_{i} + \delta_{q_{i-1}}(X_i) - \delta_{q_i}\left(\coprod_{j=0}^i \gf_{q_j}(R'_j)_{1: N}\right)\right)\\
% 		&\overset{(\ref{eq:generalized B_n as concatenation of greedy Z i})}{=}  \chi_i^N \left(\frac{1}{2^{\Sigma_i}} S'_{i} + \delta_{q_{i-1}}(X_i) - \delta_{q_i}\left(X_{i+1}\right)\right), \label{eq:X i + 1 X i relationship}
% 	\end{align}
% and
% \begin{align}
% 	R_1 - q_0^NR_0 &= q_0^N(1-\xi_0^N)R_0 + \xi_0^N G_{q_0}^N(R'_0) - q_0^N R_0 = \xi_0^N G_{q_0}^N(R'_0) = G_{q_0}^N(R'_0) \\
% 	& \overset{(\ref{eq:n iterations of greedy map is same as reading the first n bits})}{=} q_0^N\left( R'_0 - \delta_{q_0}(\gf_{q_0}(R'_0)_{1: N}) \right) \overset{(\ref{eq:generalized B_n as concatenation of greedy Z i})}{=} q_0^N\left( \frac{1}{2^{\Sigma_0}}S'_0 - \delta_{q_0}(B_1) \right)\label{eq:X 1 X 0 relationship}.
% \end{align}
% Moreover, note that
% \begin{align}
% 	& \sum_{i=0}^n \left( \prod_{j=i+1}^{n+1} q_j^N\right) \left(R_{i+1} - q_i^N R_i \right) =\sum_{i=0}^n \left(R_{i+1}\prod_{j=i+1}^{n+1} q_j^N - R_i\prod_{j=i}^{n+1} q_j^N \right)\\
% 	&=\sum_{i=1}^{n+1} \left(R_i\prod_{j=i}^{n+1} q_j^N\right) -\sum_{i=0}^n \left( R_i\prod_{j=i}^{n+1} q_j^N \right) 
% 	= R_{n+1} \prod_{j=n+1}^{n+1} q_j^N - R_0 \prod_{j=0}^{n+1} q_j^N\\
% 	&=R_{n+1}q_{n+1}^N - R_0 \chi_{n+1} \overset{(\ref{eq:definition of the generalized dynamical system})}{=} R_{n+1}q_{n+1}^N. \label{eq:relation X n + 1 with sum of X n s}
% \end{align}
% Then, 
% \begin{align}
% 	R_{n+1} &\overset{(\ref{eq:relation X n + 1 with sum of X n s})}{=} \frac{1}{q_{n+1}^N} \sum_{i=0}^n \left( \prod_{j=i+1}^{n+1} q_j^N\right) \left(R_{i+1} - q_i^N R_i \right) =  \sum_{i=0}^n \left( \prod_{j=i+1}^{n} q_j^N\right) \left(R_{i+1} - q_i^N R_i \right)\\
% 	& = \sum_{i=1}^n \left( \prod_{j=i+1}^{n} q_j^N\right) \left(R_{i+1} - q_i^N R_i \right) + \left( \prod_{j=1}^{n} q_j^N\right) \left(R_{1} - q_0^N R_0 \right)\\
% 	&\overset{(\ref{eq:X i + 1 X i relationship})}{=} \sum_{i=1}^n \left( \prod_{j=i+1}^{n} q_j^N\right) \chi_i^N \left(\frac{1}{2^{\Sigma_i}} S'_{i} + \delta_{q_{i-1}}(X_i) - \delta_{q_i}\left(X_{i+1}\right)\right) + \left( \prod_{j=1}^{n} q_j^N\right) \left(R_{1} - q_0^N R_0 \right) \nonumber\\
% 	&\overset{(\ref{eq:X 1 X 0 relationship})}{=} \sum_{i=1}^n  \chi_n^N \left(\frac{1}{2^{\Sigma_i}} S'_{i} + \delta_{q_{i-1}}(X_i) - \delta_{q_i}\left(X_{i+1}\right)\right) + \left( \prod_{j=1}^{n} q_j^N\right) q_0^N\left( \frac{1}{2^{\Sigma_0}}S'_0 - \delta_{q_0}(B_1)\right) \nonumber\\
% 	&= \chi_n^N  \left(\sum_{i=0}^n \frac{1}{2^{\Sigma_i}} S'_{i} + \sum_{i=1}^n \delta_{q_{i-1}}(X_i) - \sum_{i=0}^n\delta_{q_i}\left(X_{i+1}\right)\right)\\
% 	&= \chi_n^N  \left(\sum_{i=0}^n \frac{1}{2^{\Sigma_i}} S'_{i} - \delta_{q_n}\left(X_{n+1}\right)\right). \label{eq:intermediate result X n +1}
% \end{align}
% Remark that, for $i \in \{0,\ldots,n\}$,
% \begin{align}
% 	S'_{i} &=  S_{i} - G_2^{M_{i}}(S_{i}) \overset{(\ref{eq:n iterations of greedy map is same as reading the first n bits})}{=} \delta_2\left(\gf_2(S_{i})\right) 
% 	= \delta_2\left(\gf_2\left(G_2^{\Sigma_{i}}(S_0) \right)_{1: M_{i}}\right)\\
% 	&\overset{(\ref{eq:n iterations of greedy map is same as shifting n bits})}{=} \delta_2\left(\gf_2(S_0)_{\Sigma_{i}+1:\Sigma_{i}+M_{i}}\right)
% 	 =\delta_2\left(\gf_2(S_0)_{\Sigma_{i}+1:\Sigma_{i+1}}\right)
% 	 \overset{(\ref{eq:definition of the dynamical system})}{=}\delta_2\left(\gf_2(s)_{\Sigma_{i}+1:\Sigma_{i+1}}\right). \label{eq:generalized W_i equivalent to delta of greedy binary}
% \end{align} 
% Therefore,
% \begin{equation}
% 	\sum_{i=0}^n \frac{1}{2^{\Sigma_i}} S'_i \overset{(\ref{eq:generalized W_i equivalent to delta of greedy binary})}{=} \sum_{i=0}^n \frac{1}{2^{\Sigma_i}} \delta_2\left(\gf_2(s)_{\Sigma_{i}+1:\Sigma_{i+1}}\right) \overset{\ref{lem:factorisation of beta expansion through concatenation}}{=} \delta_2\left(\gf_2(s)_{1: \Sigma_{i+1}}\right).\label{eq:sum W i to make x}
% \end{equation}
% Finally,
% \begin{equation}
% 	\delta_{q_n}\left(X_{n+1}\right) \overset{(\ref{eq:intermediate result X n +1})}{=} \sum_{i=0}^n \frac{1}{2^{\Sigma_i}} S'_{i}  - \frac{1}{\chi_n^N}X_{n+1} \overset{(\ref{eq:sum W i to make x})}{=} \delta_2\left(\gf_2(s)_{1: \Sigma_{n+1}}\right) - \frac{1}{\chi_n^N}R_{n+1}.
% \end{equation}
% As $\DD_\beta$ satisfies $(\mathbf{H}')$, by Lemma \ref{lem:generalized d beta is well defined}, $0 \leq R_{n+1} \leq \frac{2}{\beta}$, and as $\chi_n \overset{(\ref{eq:def of chi i})} \geq q_0^n$, then $0 \leq \limi{n} \frac{R_{n+1}}{\chi_n} \leq \limi{n} \frac{2}{\beta q_0^n} = 0$. Moreover, as $\DD_\beta$ satisfies $(\mathbf{H}')$, then also $\Sigma_n \geq N n \log_2(\beta)$, so $\limi{n}\Sigma_n = \infty$. We deduce that
% \begin{equation}
% 	\limi{n} \delta_{q_n}(X_{n+1}) = \limi{n} \delta_2\left(\gf_2(s)_{1: \Sigma_{n+1}}\right) - \limi{n}\frac{1}{\chi_n^N}R_{n+1} = \delta_2(\gf_2(s)) = s,
% \end{equation}
% thereby finishing the proof.
% \end{proof}

\subsection{\texorpdfstring{Algorithmic complexity of the $\beta$-expansion generated by the dynamical system}{}}

We have established that for $\beta \in (1,2)$, and $s \in [0,1]$, $\EE_N(\ssf,\betaf)$ constitutes a valid $\beta$-expansion, assuming that (\ref{eq:final value parameters beta r 1}), (\ref{eq:beta i to generate beta expansion beta r}), (\ref{eq:final value parameters beta r 2}) and  (\ref{eq:s i to generate beta expansion beta r}) hold. In this paragraph, we conclude the proof of (\ref{eq:inequality to be proven}), which establishes a relationship between the algorithmic complexity of the respective prefixes of $\EE_N(\ssf,\betaf)$ and $\gf_2(s)$, respectively.

\begin{lemma}\label{lem:weak inequality binary beta expansion Kolmogorov rational beta r}
	Let $\beta \in (1,2)$ and $s \in [0,1]$. Define $N$, $L$ and $(\Lambda_i)_{i \in \N_0}$ according to (\ref{eq:final value parameters beta r 1}), $\betaf = \beta_0\beta_1\ldots$ according to (\ref{eq:beta i to generate beta expansion beta r}), $(\Sigma_i)_{i \in \N_0}$  and (\ref{eq:final value parameters beta r 2}),  and $\ssf = s_0s_1\ldots $ according to (\ref{eq:s i to generate beta expansion beta r}). Then, there exists $\xf \in \Sigma_\beta(s)$, such that
	\begin{equation}
		K[\xf|N n] \leq K[\gf_2(s)|\Sigma_n] +F(K[\gf_2(\betaun)|Nn+L])+ c_\beta, \ \ \text{for all} \ n \in \N.
	\end{equation}
	where $F : t \mapsto t + 2\log_2(t)$, and $c_\beta > 0$ depends only on $\beta$.
\end{lemma}
\begin{proof}
	Let $\beta \in (1,2)$, $s \in [0,1]$. Define $N$, $L$ and $(\Lambda_i)_{i \in \N_0}$ according to (\ref{eq:final value parameters beta r 1}), $\betaf = \beta_0\beta_1\ldots$ according to (\ref{eq:beta i to generate beta expansion beta r}), $(\Sigma_i)_{i \in \N_0}$  and (\ref{eq:final value parameters beta r 2}),  and $\ssf = s_0s_1\ldots $ according to (\ref{eq:s i to generate beta expansion beta r}). By Theorem \ref{thm:convergence rational dynamical system to beta expansion beta r}, $\EE_N(s_0s_1\ldots s_{n-1})$ is an $Nn$-prefix of $\EE_N(\ssf) \in \Sigma_\beta(s)$, for $n \in \N$. The equations of (\ref{eq:R i X i + 1 definition of E N}) can be run on a digital computer because they require only finite precision, and hence $\EE_N$ can be cast into the effective algorithm $\varphi : \{0,1\}^\ast \to \{0,1\}^\ast$ summarized in Algorithm \ref{alg:conversion binary expansion to beta expansion beta r}. Note that $\varphi(\langle\gf_2(s)_{1:\Sigma_n}, \gf_2(\beta-1)_{1:Nn+L}\rangle) = \EE_N(s_0s_1\ldots s_{n-1},\beta_0\beta_1\ldots \beta_{n-1}) = \EE_N(\ssf,\betaf)_{1:Nn}$, for $n \in \N$. Therefore, by Lemma \ref{lem:if two sequences are connected through an algorithm, they have same kolmogorov complexity} there exists $c_\varphi > 0$ such that
	\begin{align}
		K[\xf|N n] &\leq K[\langle\gf_2(s)_{1:\Sigma_n}, \gf_2(\beta-1)_{1:Nn+L}\rangle] + c_\varphi,\\
		&\overset{(\ref{eq:simple inequality on pairing of sequences})}{\leq}  K[\gf_2(s)|\Sigma_n] +F(K[\gf_2(\beta-1)|Nn+L]) + C_+ + c_\varphi,
	\end{align}
	for all $n \in \N$. $\varphi$ depends only on $\EE_N$, which itself depends only on $\beta$ as $N = \lceil \log_{\betaprox}(2) \rceil$. Hence, $c_\varphi$ depends only on $\beta$.
	\begin{algorithm}[H]
		\caption{$\beta$-expansion from a binary expansion}\label{alg:conversion binary expansion to beta expansion beta r}
		\begin{algorithmic}[1]
			\Require $n \in \N$, $x \in \{0,1\}^\ast$, $y \in \{0,1\}^\ast$
			\State Check that $|y| \geq Nn+L$, if not raise an error.
			\For{$i = 0, \ldots, n$}
			\State $\beta_i \gets 1 + \sum_{j=1}^{Ni+L} y_{j} 2^{-j}$
			\EndFor
			\For{$i = 0, \ldots, n$}
			\State $\Sigma_i \gets \lceil Ni\log_2(\beta_{i+1}) \rceil - \lfloor \log_2(C_\beta) \rfloor$
			\EndFor
			\State Check that $|x| \geq \Sigma_n$, if not raise an error.
			\For{$i = 0, \ldots, n-1$}
			\State $s_i \gets \frac{\beta_{i+1}^{Ni}}{2^{\Sigma_i}}\delta_2(x_{\Sigma_i+1:\Sigma_{i+1}})$
			\EndFor
			\State \Return $\EE_N(s_0s_1\ldots s_{n-1}, \beta_0\beta_1)$
		\end{algorithmic}
	\end{algorithm}
\end{proof}

% Having established that $\beta \in (1,2)$ and $s \in [0,1]$, $\xf_\beta(s)$ constitutes a valid $\beta$-expansion under condition $(\mathbf{H}')$, it remains to show that the algorithmic complexity of the prefixes of $\xf_\beta(s)$ can inferred from the algorithmic complexity of the prefixes of the greedy expansion $\gf_2(s)$. In this section, we say that $\DD_\beta$ is effective if there exists an effective algorithm which can, on input $n$, output $\Sigma_n$, for all $n \in \N$. We start with an intermediate result, and then conclude on the final result as a corollary. 
% \begin{theorem}\label{lem:generalized weak inequality binary beta expansion Kolmogorov}
% 	Let $\beta \in (1,2)$, $s \in [0,1]$, $\DD_\beta = (L,N,(\Sigma_n))$ be effective, and suppose that there exists $\delta \in \left(0, \frac{2 -\beta}{\beta}\right) \cap \Q$ such that $\DD_\beta$ satisfies \textnormal{(\ref{eq:H delta prime})}. Then, there exists $c > 0$, depending only on $\DD_\beta$, such that
% 	\begin{equation}
% 		K[\xf_\beta(s)|N n] \leq K[\gf_2(s)|\Sigma_n] +F(K[\gf_2(\beta-1)|Nn+L]) + c, \ \ \text{for all} \ n \in \N,
% 	\end{equation}
% 	where $F : t \mapsto t + 2\log_2(t)$.
% \end{theorem}
% \begin{proof}
% 	Let $\beta \in (1,2)$, $s \in [0,1]$, $\DD_\beta = (L,N,(\Sigma_n))$ be effective and suppose that there exists $\delta \in \left(0, \frac{2 -\beta}{\beta}\right) \cap \Q$ such that $\DD_\beta$ satisfies (\ref{eq:H delta prime}).
	
% 	We design an effective algorithm $\varphi : \{0,1\}^\ast \to \{0,1\}^\ast$ such that
% \begin{equation}
% 	\xf_\beta(s)_{1: N n} = \varphi(\langle \gf_2(s)_{1: \Sigma_n}, \gf_2(\beta-1)_{1: Nn+L}  \rangle), \ \ \text{for all} \ n \in \N.
% \end{equation}
% Let $n \in \N$. Concretely, on input $\langle \gf_2(\beta-1)_{1: Nn+L},  \gf_2(s)_{1: \Sigma_n} \rangle$, $\varphi$ first separates $\gf_2(\beta-1)_{1:  Nn+L}$ from $\gf_2(x)_{1: \Sigma_n}$. $\varphi$ then computes $n$ through $n = \frac{Nn+L - L}{N}$. Further, $\varphi$ uses $\gf_2(\beta-1)_{1:  Nn+L}$ to compute $q_0,\ldots, q_n$, as per $q_i = \delta_2(\gf_2(\beta-1)_{1:  Ni+L})$.
% $\varphi$ then computes $n$ iterations of the dynamical system $\DD_\beta$ on input $s$, and records the successive values $R'_0, \ldots, R'_n$. Finally, it generates $\xf_\beta(s)_{1: N n}$ through the equation
% \begin{equation}
% 	\xf_\beta(s)_{1: N n} :=  \coprod_{i=0}^{n-1} \gf_{q_i}(R'_i)_{1: N}.
% \end{equation}
% Note that as $q_k \in \Q$ for all $0\leq k \leq n$, all the variables involved in the evolution of $\DD_\beta$ are rational numbers, and all the operations underlying this evolution can be run on a digital computer, as they can be reduced to additions and multiplications only. So $\varphi$ is indeed an effective algorithm. By Lemma \ref{lem:if two sequences are connected through an algorithm, they have same kolmogorov complexity}, there exists $C > 0$, depending only on $\varphi$, and a fortirori only on $\DD_\beta$, such that, for all $n \in \N$,
% \begin{align}
% 	K[\xf_\beta(x)|N n] &\leq K[\langle \gf_2(s)_{1: \Sigma_n}, \gf_2(\beta-1)_{1: Nn+L}  \rangle] + C\\
% 	&\overset{(\ref{eq:simple inequality on pairing of sequences})}{\leq} K[\gf_2(s)_{1: \Sigma_n}] + F(K[\gf_2(\beta-1)_{1:  Nn+L}]) + C_+ + C.
% \end{align}
% We set $c := C + C_+$, thereby finishing the proof.
% \end{proof}

\begin{corollary}
	Let $\beta \in (1,2)$, $s \in [0,1]$. Then, there exists $\xf\in \Sigma_\beta(s)$ such that
	\begin{equation}\label{eq:generalized algorithmic complexity of beta expansion lower than binary expansion}
		K[\xf|n] \leq  K[\gf_2(s)|\lceil n \log_2(\beta)\rceil] + F(K[\gf_2(\betaun)|n]) + c_\beta, \ \text{for all} \ n \in \N.
	\end{equation}
	where $F : t \mapsto t + 2\log_2(t)$, and $c_\beta > 0$ depends only on $\beta$.
\end{corollary}
\begin{proof}
	\textit{To appear.}
\end{proof}
% \begin{proof}
% 	Let $\beta \in (1,2)$, $s \in [0,1]$. Let $\delta \in (0,\frac{2-\beta}{\beta}) \cap \Q$, and let $\DD_\beta$ such that 
% 	\begin{equation}\label{eq:final conditions of N and Sigma, H delta prime in proof}
% 			\tag{$\mathbf{H}_\delta'$}
% 			\begin{split}
% 			L &= \left\lceil\max\left\{C_{N,\beta} -\log_2\left(\frac{C_{\delta,\beta}}{N^2\beta^N}\right) , \frac{1}{2} \log_2\left(-\log\left(C'_{N,\delta,\beta}\right)\right)\right\}\right\rceil\\
% 			N &= -  \left\lceil\log_\beta\left( \frac{\delta}{2- \beta} \right)\right\rceil, \\
% 			\Sigma_{i} &= \left\lceil Ni\log_2(q_i) - \log_2\left(C_{\delta,\beta} \right)\right\rceil , \ \forall i \in \N.
% 			\end{split}
% 	\end{equation}
% 	Obviously, $\DD_\beta$ satisfies (\ref{eq:H delta prime}), where $\delta \in \Q$, so by Lemma \ref{lem:generalized weak inequality binary beta expansion Kolmogorov}, there exists $c>0$ such that
% 	\begin{equation}
% 		K[\xf_\beta(s)|N n] \leq K[\gf_2(s)|\Sigma_n] +F(K[\gf_2(\beta-1)|Nn+L]) + c, \ \ \text{for all} \ n \in \N.
% 	\end{equation}
% 	Note that by (\ref{eq:simple inequality on concatenation of sequences}), there exists $C>0$ such that
% \begin{equation}\label{eq:generalized before ultimate equation proof kolmogorov}
% 	K[\xf_\beta(s)|n]  \leq K[\gf_2(s)|\Sigma_{\lfloor n/N\rfloor}]+F(K[\gf_2(\beta-1)|N\lfloor n/N\rfloor+L]) + c + C, \ \ \text{for all} \ n \in \N.
% \end{equation}
% Further, remark that
% \begin{equation}
% 	n - N + L \leq N\lfloor n/N\rfloor+L \leq n + L, \ \text{ for all } \ n \in \N.
% \end{equation}
% Then, by Lemma \ref{lem:kolmogorov complexity of closed sequences}, there exists $C' > 0$ such that
% \begin{align}
% 	& K[\gf_2(\beta-1)|N\lfloor n/N\rfloor+L] \leq K[\gf_2(\beta-1)|n] + C', \ \ \text{for all} \ n \in \N,\\
% 	\overset{(a)}{\Rightarrow} & F(K[\gf_2(\beta-1)|N\lfloor n/N\rfloor+L]) \leq F(K[\gf_2(\beta-1)|n]) + C', \ \ \text{for all} \ n \in \N,
% \end{align}
% where (a) is a consequence of $F$ being increasing and subadditive. It follows that 
% \begin{equation}\label{eq:generalized before ultimate equation proof kolmogorov 2}
% 	K[\xf_\beta(s)|n]  \leq K[\gf_2(s)|\Sigma_{\lfloor n/N\rfloor}]+F(K[\gf_2(\beta-1)|n]) + C' + c + C, \ \ \text{for all} \ n \in \N.
% \end{equation}
% Finally, note that
% \begin{align}
% 	\left|\lceil n \log_2(\beta)\rceil - \Sigma_{\lfloor n/N\rfloor}\right| &\leq 2 + \left|n \log_2(\beta) - N\left\lfloor \frac{n}{N}\right\rfloor\log_2(q_{\lfloor n/N \rfloor}) - \log_2\left(C_{\delta,\beta} \right)\right|\\
% 	&\leq 2 + N + \left|n \log_2(\beta) - n\log_2(q_{\lfloor n/N \rfloor}) - \log_2\left(C_{\delta,\beta} \right)\right|\\
% 	&\overset{(a)}{\leq} 2 + N + \left|n\frac{1}{\log(2)} (\beta - q_{\lfloor n/N \rfloor}) - \log_2\left(C_{\delta,\beta} \right)\right|\\
% 	&\leq 2 + N + \left|n\frac{1}{\log(2)} 2^{-L-N{\lfloor n/N \rfloor}} - \log_2\left(C_{\delta,\beta} \right)\right|\\
% 	&\leq 2 + N + \left|\frac{2}{\log(2)} 2^{-L}n2^{-n} - \log_2\left(C_{\delta,\beta} \right)\right|\\
% 	&\overset{(b)}{\leq} 2 + N + \left|\frac{2}{\log(2)} 2^{-L} - \log_2\left(C_{\delta,\beta} \right)\right| =: M\\
% \end{align}
% where (a) follows from $\log : [1, \infty] \to \R$ being 1-Lipschitz, and (b) is by $n2^{-n}\leq 1$ for all $n \in \N$. Then, $\left|\lceil n \log_2(\beta)\rceil - \Sigma_{\lfloor n/N\rfloor}\right| \leq M$ for all $n \in \N$, so by Lemma by Lemma \ref{lem:kolmogorov complexity of closed sequences}, there exists $C'' > 0$ such that
% \begin{equation}\label{eq:generalized before ultimate equation proof kolmogorov 3}
% 	K[\gf_2(s)|\Sigma_{\lfloor n/N\rfloor}] \leq K[\gf_2(s)|\lceil n \log_2(\beta) \rceil] + C'', \ \text{for all} \ n \in \N.
% \end{equation}
% By combining (\ref{eq:before ultimate equation proof kolmogorov}) and (\ref{eq:before ultimate equation proof kolmogorov 2}), one has
% \begin{equation}
% \nonumber
% 	K[\xf_\beta(s)|n] \leq K[\gf_2(s)|\lceil n \log_2(\beta) \rceil] + F(K[\gf_2(\beta-1)|n]) + C'' + C' + C + c, \ \text{for all} \ n \in \N,
% \end{equation}
% thereby finishing the proof by letting $\hat \xf_\beta := \xf_\beta (s)$.
% \end{proof}

We have proven that for every $s \in [0,1]$, there exists a $\beta$-expansion of $s$ which algorithmic complexity of prefixes is upper bounded a very simple manner by a relation involving the algorithmic complexity of the prefixes of the binary expansion of $s$ and the binary expansion of $\beta$, for every $\beta \in (1,2)$. A direct consequence is the following corollary, when $\beta$ is a computable number.
\begin{corollary}
	Let $\beta \in (1,2)$ be a computable number, and $s \in [0,1]$. Then, there exists $\xf \in \Sigma_\beta(s)$ such that
	\begin{equation}\label{eq:computable algorithmic complexity of beta expansion lower than binary expansion}
		K[\xf|n] \leq  K[\gf_2(s)|\lceil n \log_2(\beta)\rceil] + c_\beta, \ \text{for all} \ n \in \N,
	\end{equation}
	where $c_\beta > 0$ depends only on $\beta$.
\end{corollary}
\noindent This achieves the proof of the main Theorem.
